\subsection{Benchmark Models}\label{sec:benchmarks}
The generated option prices are evaluated against two benchmark models, a GBM and a historical simulation of the realized returns. Both benchmark models replaced $G$ with a simpler model of the returns and were priced through the procedure shown in Fig.~\ref{fig:pricing-pipeline} without change. The benchmarks were run with a single seed each and the results do not carry a seed spread.

GBM is the process underlying Black-76, introduced in Sec.~\ref{sec:implied-vol}. Under GBM the return over a short period is normally distributed and returns over different periods are independent. The price at a future date is lognormally distributed \parencite[332]{hullOptionsFuturesOther2022}. The futures price is not expected to grow, so the log return from the valuation date to expiry is normally distributed
\begin{equation}
    \ln \frac{F_T}{F_t} \sim
    \mathcal{N}\!\left(-\tfrac{1}{2}\sigma^{2}T, \; \sigma^{2}T\right),
    \label{eq:gbm-horizon}
\end{equation}
where $-\frac{1}{2}\sigma^2 T$ and $\sigma^2 T$ are the mean and variance of the log return. $-\frac{1}{2}\sigma^2 T$ follows from the risk-neutral condition $\hat{\mathbb{E}}[F_T] = F_t$ of Sec.~\ref{sec:options} \parencite[328]{hullOptionsFuturesOther2022}. Each path of the GBM was simulated as $h$ returns over the trading days to expiry, 
\begin{equation}
    \hat{x}_k = -\tfrac{1}{2}\sigma^{2}\Delta
          + \sigma\sqrt{\Delta}\,\varepsilon_k,
    \qquad
    \varepsilon_k \sim \mathcal{N}(0,1),
    \label{eq:gbm-discrete}
\end{equation}
where $\Delta = T / h$ is the time step and $\varepsilon_k$ are independent standard normal random variables. The returns are drawn independently, so the variance of the sum of $h$ returns is $h \sigma^2 \Delta = \sigma^2 T$.

GBM reproduces none of the stylized facts of Sec.~\ref{sec:stylized-facts}. The returns are normally distributed, so the simulated distribution does not reproduce heavy tails or asymmetry. The returns are also drawn independently, so there is no volatility clustering. A single $\sigma$ applies to every strike, so the volatility smile of the simulated prices is flat.

Two choices of $\sigma$ are considered. The first is the implied volatility published by the exchange at the listed strike closest to the settlement price on the valuation date. The published $\sigma$ is obtained by inverting Black-76 at that strike, so pricing the same strike returns the market price. The second choice is the annualized realized volatility of the front-month returns over the preceding \num{60} trading days. The realized volatility is computed from the return series alone, without option prices. The realized volatility lies below the implied volatility on \num{16} of the \num{20} valuation dates. The results of Sec.~\ref{sec:option-prices} are reported with the implied volatility.

The historical simulation draws paths of $h$ returns by resampling the observed returns, a procedure known as the bootstrap \parencite{efronBootstrapMethodsAnother1979}. The returns are drawn from the training split of Sec.~\ref{sec:data-pipeline}, so no test data enters the benchmark. Two modes of the bootstrap are considered. In the independent and identically distributed ()$\mathrm{iid}$) mode, each of the $h$ returns is drawn independently with replacement, so the same observed return can enter a path more than once. No distributional assumption is made, since the returns are drawn from the observed distribution. In the block mode, each path is a block of $h$ consecutive returns \parencite{kunschJackknifeBootstrapGeneral1989}. The order of the observed series is preserved within each block, so the volatility clustering of Sec.~\ref{sec:stylized-facts} is present as well. The training split contains \num{3178}\footnote{A block of $h$ consecutive returns can start at $3197 - h + 1$ positions, where \num{3197} is the number of returns in the training split.} distinct blocks, so the \num{20000} paths reuse each block multiple times. The paths are not independent, so the Monte Carlo standard error of Sec.~\ref{sec:pricing-setup} is too small in the block mode. 

Table~\ref{tab:benchmark-properties} summarizes the stylized facts reproduced by each model. $G$ is the only model conditioned on the returns preceding the valuation date. GBM and the $\mathrm{iid}$ bootstrap differ in the distribution of the returns, while both draw the returns independently. The difference between these two benchmarks is attributable to the heavy tails and the asymmetry of Sec.~\ref{sec:stylized-facts} jointly. The two bootstrap modes draw from the same returns and differ only in the order, so the difference between them is attributable to the volatility clustering.

\begin{table}[htbp]
  \centering
  \begin{tabular}{ll}
    \toprule
    Model & Stylized facts \\
    \midrule
    GBM & none \\
    Bootstrap, $\mathrm{iid}$            & heavy tails, asymmetry \\
    Bootstrap, block          & heavy tails, asymmetry, volatility clustering \\
    $G$                       & heavy tails, asymmetry, volatility clustering \\
    \bottomrule
  \end{tabular}
  \caption{Stylized facts reproduced by the benchmarks and $G$}
  \label{tab:benchmark-properties}
\end{table}

Further benchmark models are considered and discussed in Chapter~\ref{sec:discussion}. The prices of the benchmarks and of $G$ are reported and compared in Sec.~\ref{sec:option-prices}. The benchmarks are implemented in \texttt{price\_gbm.py} and \texttt{price\_bootstrap.py}.
